\documentclass[11pt,letterpaper]{article}
\usepackage[margin=1in]{geometry}
\usepackage{booktabs}
\usepackage{hyperref}
\usepackage{xcolor}
\usepackage{longtable}
\usepackage{parskip}
\hypersetup{colorlinks=true,linkcolor=blue,urlcolor=blue,citecolor=blue}

\title{Independent Replication of Stover \emph{et al.}\ 2000 \\
       \large Complete Genome Sequence of \emph{Pseudomonas aeruginosa} PAO1}
\author{BVBRC-100 Replication Wave \\ Target: BVBRC-98-Paeruginosa-PAO1-Stover2000}
\date{2026-07-04 \quad Host: CherryRd \quad Verdict: \textbf{PARTIAL}}

\begin{document}
\maketitle

\section{Paper under replication}

Stover C.K., Pham X.Q., Erwin A.L., \emph{et al.}
\emph{Complete genome sequence of Pseudomonas aeruginosa PAO1, an
opportunistic pathogen.} \textbf{Nature} 406(6799):959--964
(31 Aug 2000). \texttt{doi:10.1038/35023079}. PMID~10984043.

The paper reports the finished 6.26~Mbp single-circular-chromosome
sequence of strain PAO1, at the time the largest bacterial genome
sequenced, with an initial annotation of $\sim$5{,}570 predicted ORFs and
66.6\,\% G+C, and highlights the exceptional richness of regulatory and
two-component signalling genes. The sequence deposited then evolved into
GenBank/RefSeq record \texttt{NC\_002516.2} (assembly ASM676v1,
\texttt{GCF\_000006765.1}), which is the version used for this
replication.

\section{Claims tested}

\begin{longtable}{@{}p{0.05\linewidth}p{0.55\linewidth}p{0.14\linewidth}p{0.18\linewidth}@{}}
\toprule
ID & Claim & Type & Testable? \\
\midrule
C1 & Genome size $= 6{,}264{,}403$~bp, single circular chromosome & quantitative & yes \\
C2 & G+C content $= 66.6$\,\% & quantitative & yes \\
C3 & Predicted ORFs $= 5{,}570$ protein-coding genes & quantitative & yes \\
C4 & Largest sequenced bacterial genome at publication & historical/context & not from FASTA \\
C5 & Highest fraction of regulatory / two-component genes among then-sequenced bacteria & comparative/context & not from single FASTA \\
\bottomrule
\end{longtable}

\section{Method}

Environment: macOS host \texttt{CherryRd}, NCBI \texttt{datasets} CLI
v18.x, Python 3.13 stdlib only. Steps:

\begin{enumerate}
\item Fetch RefSeq assembly:
\begin{verbatim}
datasets download genome accession GCF_000006765.1 \
    --include genome,gff3,protein
unzip ncbi_dataset.zip
\end{verbatim}
Yields \texttt{GCF\_000006765.1\_ASM676v1\_genomic.fna} (6.3~MB),
\texttt{genomic.gff} (3.2~MB), \texttt{protein.faa} (2.3~MB); MD5s
recorded in \texttt{report/artifact\_harvest.md}.
\item Parse FASTA (\texttt{work/analyze.py}): uppercase, count A/C/G/T,
compute total length and $\mathrm{GC\%} = 100\,(G+C)/(A+C+G+T)$.
\item Parse GFF3: tally \texttt{gene}, \texttt{CDS}, \texttt{rRNA},
\texttt{tRNA}, \texttt{ncRNA}, \texttt{tmRNA}; extract
\texttt{gene\_biotype} + \texttt{protein\_id}; count unique protein IDs.
\item Sanity-check \texttt{protein.faa}: count \texttt{>} headers vs.\
unique protein IDs from GFF.
\item Provenance: MD5 every downloaded file.
\item LLM-judge (\texttt{work/llm\_judge.py}) via free Argo endpoint
\texttt{http://127.0.0.1:44497/v1/chat/completions} model
\texttt{argo:gpt-4o}, \texttt{temperature=0.0}. Full request/response
cached to \texttt{report/evidence/llm\_judge.json}.
\end{enumerate}

\section{Results}

\begin{longtable}{@{}p{0.06\linewidth}p{0.20\linewidth}p{0.28\linewidth}p{0.18\linewidth}p{0.16\linewidth}@{}}
\toprule
Claim & Paper value & Replication (NC\_002516.2, 2026-07-04) & $\Delta$ & Reproduced? \\
\midrule
C1 & $6{,}264{,}403$~bp & $\mathbf{6{,}264{,}404}$ & $+1$~bp ($+1.6\times10^{-5}$\,\%) & effectively exact \\
C2 & $66.6\,\%$ & $\mathbf{66.556\,\%}$ & $-0.044$~pp & within rounding \\
C3 & $5{,}570$ ORFs & $\mathbf{5{,}573}$ CDS (unique protein IDs $=5{,}572$; \texttt{protein.faa} $=5{,}572$) & $+3$ ($+0.054$\,\%) & within annotation drift \\
tRNA & $\sim$63 & $\mathbf{63}$ & --- & consistent \\
rRNA operons & 4 & $13$ rRNA features $/\,3=4$ & --- & consistent \\
topology & single circular & 1 contig, 0 ambiguous bases & --- & consistent \\
coding density & $\sim$89\,\% & $\mathbf{89.3\,\%}$ & --- & consistent \\
C4 & true & not re-derivable from FASTA; historically corroborated (e.g.\ Nierman 2001 \emph{B.~pseudomallei} 7.2~Mbp, one year later) & --- & context-only \\
C5 & true & not re-derivable from single FASTA; corroborated (Rodrigue 2000; Galperin 2005; PMC9607943 2022) & --- & context-only \\
\bottomrule
\end{longtable}

LLM-judge (\texttt{argo:gpt-4o}, $T=0$): overall verdict \textbf{PARTIAL}
(C1~=~yes; C2/C3~=~``partial'' because sub-0.1\,\% deviations flagged;
C4/C5~=~not-testable-from-genome). JSON cached to
\texttt{report/evidence/llm\_judge.json}.

\section{Verdict}

\textbf{PARTIAL}, adopting the LLM-judge verdict as canonical per the
wave-brief rule. All three numerically-testable claims (genome size, GC
content, ORF count) reproduce with essentially exact agreement ($\Delta
= +1$~bp, $-0.044$~pp, $+3$~CDS respectively --- cumulatively well under
the noise floor for a 25-year-old annotation). The two remaining paper
claims (C4~largest-bacterial-genome-at-publication, C5~exceptional
regulatory-gene richness) are historical/comparative statements about
the year-2000 sequenced-genome landscape and cannot be re-derived from a
single FASTA; both are uncontested in follow-up literature. Under a
``context/spot-check'' scoring rubric this replication would be a clean
REPLICATED; it is recorded as PARTIAL to stay conservative and honor the
judge's verdict.

\section{Genuine critique}

\paragraph{What genuinely replicated.} The core quantitative backbone of
the Stover \emph{et al.}\ 2000 paper --- the three numbers most cited
downstream (6.26~Mbp, 66.6\,\% G+C, $\sim$5{,}570 ORFs) --- reproduces
to within a single base pair, $\sim$0.04 percentage points, and three
protein-coding features. That level of stability across a 25-year gap
is genuinely impressive and reflects (i) the high finishing quality of
the original Sanger-era submission and (ii) the fact that RefSeq
\texttt{NC\_002516.2} is a lightly-revised descendant of the same
submitted sequence, not an independently-resequenced assembly. The
tRNA count, rRNA-operon count, single-circular-chromosome topology, and
$\sim$89\,\% coding density all match the paper's narrative claims as
well. If the assignment were purely ``do the numbers survive?'', this
would be a clean pass.

\paragraph{What did \emph{not} replicate --- and why the verdict is
PARTIAL, not REPLICATED.} Two of the five paper claims (C4, C5) are
comparative/historical statements about the year-2000 sequenced-genome
landscape: \emph{largest bacterial genome sequenced to that date}, and
\emph{highest fraction of regulatory/two-component genes among
then-sequenced bacteria}. Neither of these can be tested against a
single downloaded FASTA. They would require rebuilding a year-2000
sequenced-genome baseline (which the replication did not do) and
recomputing regulatory-gene fractions across that baseline under a
harmonised annotation pipeline (also not done). The follow-up literature
is uncontested on both points, but ``uncontested in the literature'' is
not the same as ``independently replicated here,'' and the honest score
for those two claims in this exercise is \emph{not tested}, not
\emph{reproduced}. The LLM-judge (\texttt{argo:gpt-4o}, $T=0$) also
downgraded C2/C3 to ``partial'' on the strength of the sub-0.1\,\%
numerical deviations; a strict reading of that judgement is defensible
even if practically the drift is trivial.

\paragraph{What this replication is not.} This is a numerical
spot-check against the current RefSeq record, not a re-assembly from
raw reads (the Sanger reads are effectively not re-runnable), not an
independent re-annotation (we consumed the PGAP-produced GFF instead of
running a fresh Prokka/Bakta pass), and not a comparative-genomics
regeneration of C5. A stronger replication would (a) run an independent
annotation pipeline on the same FASTA and compare ORF counts head-to-head
with both the paper (5{,}570) and the current RefSeq (5{,}573), and (b)
rebuild a year-2000 bacterial-genome cohort and recompute regulatory-gene
enrichment to actually test C5. Neither was done.

\paragraph{Reference-strain caveat.} The single ``PAO1'' label in the
paper conceals real biological heterogeneity: PAO1 is maintained as
multiple sublineages in strain collections worldwide
(Nottingham/Zurich/UCSF/MPAO1/PAO1-DSM/etc.), and these sublineages
carry documented SNPs, small indels, and mobile-element differences.
The Stover 2000 sequence corresponds to one specific isolate, and
downstream functional-genomics work performed on other PAO1 sublineages
will not, in general, map 1:1 to \texttt{NC\_002516.2}. This is a
reproducibility concern for the wider PAO1 literature, not a flaw of
Stover 2000 itself, but a genuine replication effort should surface it.

\paragraph{Bottom line.} 3/5 claims reproduced to within noise; 2/5
claims not testable from the downloaded artefact. Verdict PARTIAL is the
right honest call.

\section{Constraints honored}

Free endpoints only (NCBI \texttt{datasets} + Argo local proxy, free
\texttt{argo:gpt-4o}); no Anthropic/OpenAI/OpenRouter calls; LLM-judge
supplied the final verdict; all writes stayed inside the target
directory; real replication on real bytes (6.3~MB FASTA + 3.2~MB GFF).

\end{document}
